\chapter{Исходный код и описание структуры репозитория}\label{app:A}

Исходный код разработанной системы размещён в открытом репозитории на платформе GitHub:

\begin{center}
\url{https://github.com/ButakovBI/AutoRepairEstimator}
\end{center}

Система построена на микросервисной архитектуре. Основной код системы
находится в директории src/auto\_repair\_estimator/ и разделён на три сервиса:

\begin{itemize}
    \item backend/ - FastAPI-сервис: API, бизнес-логика, компоненты отказоустойчивости;
    \item bot/ - VK-бот: интерфейс взаимодействия с пользователем;
    \item ml\_worker/ - ML-воркер: обработка изображений и инференс.
\end{itemize}

Прочие директории репозитория:

\begin{itemize}
    \item ml/ - скрипты и материалы для обучения моделей;
    \item docker/ - файлы Dockerfile и docker-compose.yml;
    \item tests/ - юнит- и интеграционные тесты;
    \item scripts/ - вспомогательные скрипты, в том числе нагрузочное тестирование;
    \item .github/workflows/ - конфигурация CI/CD пайплайна на GitHub Actions.
\end{itemize}
